\chapter{Altitude Sickness}

\section*{Definition}
Altitude sickness, also called acute mountain sickness (AMS), is a group of symptoms caused by travel to high elevation without enough acclimatization. AMS usually occurs above about 2,500 meters (8,000 feet), but susceptibility and ascent rate vary. High-altitude cerebral edema (HACE) and high-altitude pulmonary edema (HAPE) are life-threatening forms of high-altitude illness.

\section*{When to See a Doctor}
Altitude sickness requires prompt attention because it can progress quickly. Seek medical care immediately if:
\begin{itemize}
  \item You develop confusion, staggering, or difficulty walking (possible HACE — high-altitude cerebral edema)
  \item You have shortness of breath at rest, rapidly worsening breathlessness, a persistent cough, or coughing up pink or frothy sputum (possible HAPE -- high-altitude pulmonary edema)
  \item Your symptoms worsen despite resting at the same altitude
  \item You are unable to walk a straight line or have trouble with coordination
  \item You become very drowsy or difficult to wake
  \item You have severe headache, repeated vomiting, bluish lips, confusion, or severe weakness
\end{itemize}

\section*{How It Develops}
As you climb, air pressure falls and each breath provides less oxygen. The body normally adapts by breathing more, changing blood flow, and gradually increasing red-cell production. If ascent is too rapid, low oxygen can cause headache and other AMS symptoms. In severe illness, fluid accumulates in the lungs (HAPE) or brain (HACE), impairing breathing or neurologic function.

\section*{Causes}
\begin{itemize}
  \item Rapid ascent to high altitude (above 2,500 m / 8,000 ft) without gradual acclimatization
  \item Previous history of altitude sickness
  \item Physical exertion at high altitude before acclimatizing
  \item Sleeping at a higher altitude than the one you spent the day at (often described as ``climb high, sleep low'' — reversing this pattern increases risk)
  \item Individual susceptibility — some people are more prone regardless of fitness level
\end{itemize}

\section*{Related Symptoms}
\begin{itemize}
  \item Headache, the most common symptom of AMS
  \item Nausea and vomiting
  \item Dizziness or lightheadedness
  \item Fatigue and weakness
  \item Loss of appetite
  \item Shortness of breath with activity
  \item Difficulty sleeping
  \item Swelling of hands, feet, or face
\end{itemize}
\textbf{Important:} Descent is the definitive treatment for severe or worsening altitude illness. Do not leave someone with suspected HACE or HAPE alone, and do not allow an unwell person to descend without supervision and a safe evacuation plan.

\section*{Medical Evaluation}
\begin{itemize}
  \item Your doctor will ask about your recent ascent rate, the highest altitude reached, and your specific symptoms.
  \item They will check vital signs and oxygen saturation and listen to the lungs for signs of fluid. Oxygen saturation is normally lower at altitude, so it must be interpreted with the clinical picture.
  \item A neurological exam (checking coordination, mental status, and walking) helps assess for brain swelling.
  \item A chest X-ray may show fluid in the lungs if HAPE is suspected.
  \item Brain imaging (CT or MRI) may be done if HACE is suspected and the diagnosis is uncertain.
\end{itemize}

\section*{Treatment Options}
\begin{itemize}
  \item Stop ascent and descend promptly for severe, worsening, or neurologic or respiratory symptoms. Descent is definitive treatment and should not be delayed for medication.
  \item Supplemental oxygen and a portable hyperbaric chamber can stabilize a person while descent or evacuation is arranged; they do not replace descent when it is possible.
  \item Acetazolamide can speed acclimatization and treat AMS under medical guidance. It is not a substitute for descent in HACE or HAPE.
  \item Dexamethasone can reduce brain swelling in HACE and may be used for severe AMS when descent is delayed or as directed by a clinician.
  \item Nifedipine may be used for HAPE when descent or oxygen is delayed or as prevention in selected high-risk people. Phosphodiesterase-5 inhibitors are specialist options, not routine self-treatment.
  \item A portable hyperbaric chamber can buy time while awaiting evacuation.
  \item Prevention: above 3,000 m (10,000 ft), increase sleeping altitude by no more than about 500 m (1,600 ft) per day and take an extra acclimatization day for each 1,000 m gained. Consider acetazolamide prophylaxis with a clinician if risk is substantial.
\end{itemize}

\section*{Self-Care and Home Management}
\begin{itemize}
  \item Stop ascending — do not go higher until symptoms resolve completely.
  \item Rest at your current altitude and avoid strenuous activity. Do not ascend while symptoms are present.
  \item If symptoms worsen, fail to improve after a day at the same altitude, or include confusion, poor coordination, or breathlessness at rest, descend and seek emergency help immediately.
  \item Drink normally to thirst and avoid alcohol, sedative sleeping medicines, and other substances that can worsen judgment or breathing.
  \item Take acetaminophen (Tylenol) or ibuprofen (Advil, Motrin) for headache.
  \item Acetazolamide (Diamox) can help prevent and treat altitude sickness — discuss with your doctor before travel.
  \item Supplemental oxygen if available.
  \item A portable hyperbaric chamber (Gamow bag) can simulate descent in remote settings while you arrange evacuation.
\end{itemize}

\section*{References}
\begin{thebibliography}{99}
\bibitem{altitude_sickness:ref1} Luks AM, Swenson ER, et al. ``Acute high-altitude sickness.'' \textit{N Engl J Med}. 2017;377(15):1474-1483.
\bibitem{altitude_sickness:ref2} Hackett PH, Roach RC. ``High-altitude illness.'' \textit{N Engl J Med}. 2001;345(2):107-114.
\bibitem{altitude_sickness:ref3} Luks AM, Auerbach PS, et al. ``Wilderness Medical Society clinical practice guidelines for the prevention and treatment of acute altitude illness.'' \textit{Wilderness Environ Med}. 2019;30(4S):S3-S18.
\bibitem{altitude_sickness:ref4} Pollard AJ, Murdoch DR. \textit{The High Altitude Medicine Handbook}, 3rd ed. Radcliffe Publishing, 2003.
\bibitem{altitude_sickness:ref5} UpToDate. ``Acute mountain sickness, high-altitude cerebral edema, and high-altitude pulmonary edema.'' Wolters Kluwer, 2023.
\end{thebibliography}
